% RECOVERY32: reconstructed from the fully reviewed lost d9dec570 proof.
% Scope includes EVERY N=2^r n, in particular N=2n, not only odd heads.
\statement{A.101}{Proved}{Full-ring Alexander vanishing of the actual hexagon}
Let alpha be an actual negative full-GT involution, u its canonical
word, \(k=(d_\alpha-1)/2\in2\widehat{\mathbb Z}\), and
\[
 \mathcal H_d=u(x,y)x^k u(z,x)z^k u(y,z)y^k,
 \qquad z=(xy)^{-1}.
 \tag{A.101.1}
\]
For \(K_x=\overline{\langle\!\langle x\rangle\!\rangle}\),
put \(D_x=\overline{K_x^2[K_x,K_x]}\), and define \(D_y,D_z\)
by the same actual cyclic marking. Then
\[
                 \mathcal H_d\in D_x\cap D_y\cap D_z.
 \tag{A.101.2}
\]
In particular the first-cover Alexander coordinate is zero in
\(\mathbb F_2[[\widehat{\mathbb Z}]]\), including every finite
head \(N=2^r n\), n odd. This is not a full-word vanishing claim.

\begin{proof}
Use the exact A.89 module and retain
\[
\begin{gathered}
 U=[u(x,y)],\quad W=[u(y,z)],\quad B=U+\bar U,\\
 \bar W=TW,\quad \theta=1+(T-1)B,\quad
 R_d=[\mathcal H_d]=B+T^{-k}D_k+T^{-k}W.
\end{gathered}
\tag{A.101.3}
\]
All these are evaluations of the same full word with the actual
ordered powers, not independent module parameters. A.100 supplies
the full coefficient equality in the original normalization;
the argument below retains the third-pair coefficient and obtains
the stronger value of R itself.

Let R be any finite commutative ring of characteristic two, and
let \(\rho\in R^\times\) have finite order, with
\(t=1+\rho\) a unit. The entire A.100 construction and all six
action verifications hold over R: they use only multiplication
and the units rho,t. The two matrices are still involutions, and
the row vectors
\[
 \ell_+=(1,1),\qquad \ell_-=(1,\rho)
 \tag{A.101.4}
\]
are left eigenrows of \(\mathcal R=FD\) with eigenvalues
rho and \(\rho^{-1}\). Their determinant is t, a unit; on
\(w=e_2/t\) their values are \(t^{-1}\) and \(\rho/t\), both
units. Thus all finite profinite-power evaluations are legitimate,
without a field-only eigenline argument.

The first, second, third and fifth pair values from A.100 remain
valid over R by the same explicit translations. The fourth pair
is the scalar pair \((0,\rho^{-1}),(1,\rho)\) on e2. In the
original first-cover coordinates, at \(T=\rho^{-1}\) the pair
(y,z) is \((0,\rho^{-1}),(\rho,\rho)\). Scaling translations
by \(\rho^{-1}\) therefore gives its exact value
\[
 Y_\rho=\rho^{-1}W(\rho^{-1})=W(\rho),
 \tag{A.101.5}
\]
where the last equality is the actual weighted identity
\(\bar W=TW\). It is essential to retain this weight.

The same ordered pentagon gives (A.100.9) over R. Apply the first
row in (A.101.4) and cancel its unit value on w. The exact equation is
\[
 \rho^k\bigl(1+(\rho-1)B(\rho)\bigr)
                          =1+(\rho-1)W(\rho).
 \tag{A.101.6}
\]
Multiplying the formula for R in (A.101.3) by
\((\rho-1)\rho^k\) gives
\[
 (\rho-1)\rho^k R_d(\rho)
   =(\rho-1)\rho^kB(\rho)+(\rho^k-1)
                         +(\rho-1)W(\rho)=0.
\]
The geometric-sum identity used here is valid by continuity for
the actual profinite k, without dividing by a nonunit. Since
\((\rho-1)\rho^k\) is a unit in this ring, \(R_d(\rho)=0\).

Now take an arbitrary positive integer \(N=2^r n\), n odd, and
the full group ring \(R_N=\mathbb F_2[T]/(T^N-1)\). Put
\(M_n(T)=1+T+\cdots+T^{n-1}\). In characteristic two,
\[
 T^N-1=(T-1)^{2^r}M_n(T)^{2^r},\qquad M_n(1)=1.
 \tag{A.101.7}
\]
The two factors are coprime, giving the exact CRT decomposition
\[
 R_N\simeq\mathbb F_2[T]/(T-1)^{2^r}
       \times\mathbb F_2[T]/M_n(T)^{2^r}.
 \tag{A.101.8}
\]
On the second factor, when n is greater than one, T and \(T-1\) are
units, so the full ring argument just proved gives zero. This
includes its nilpotent thickening; no passage merely to residue
fields has occurred. For n=1 this factor is omitted.

The first factor is the group algebra of the cyclic group of order
\(2^r\). A.42/A.45, as explicitly used in A.89, proves that the
particular word \(\mathcal H_d\) is one in the actual maximal
pro-two quotient for the prescribed \(d_{\alpha,2}\). This is
stronger than merely knowing \(a_2=d_2\). It makes the first
factor zero too, including r=0. Thus R vanishes in every \(R_N\),
and hence in their inverse limit \(\Lambda\). This proves
\(\mathcal H_d\in D_x\).

For cyclic propagation put
\(B_1=u(x,y)x^k\), \(B_2=u(z,x)z^k\),
\(B_3=u(y,z)y^k\). The actual cyclic basis automorphism
\(\rho_c:x\mapsto y,y\mapsto z,z\mapsto x\) satisfies
\[
 \rho_c(\mathcal H_d)=B_3B_1B_2
                       =B_3\mathcal H_dB_3^{-1}.
 \tag{A.101.9}
\]
Each \(D_i\) is closed and normal in the original free profinite group,
because it is characteristic in the corresponding normal \(K_i\).
Applying the cyclic automorphism and using (A.101.9) therefore
gives membership in \(D_y\) and \(D_z\) as well. The word has zero
abelianization, so all three kernel memberships are also consistent
with their marked cyclic quotients.

These are closed square/commutator-subgroup containments, not a
claim that their intersection is trivial. In particular they do
not determine the nonlinear word or either separately normalized
CM frame. The full-GT pentagon and actual pro-two hexagon are the
precise additional inputs, and are not supplied for arbitrary GT0
or peripheral involutions.
\end{proof}
